\section{Can a Population Coordinate Without Converging?}
\label{sec:coordinate-converging}

\textbf{Unresolved issue.}
A population needs both independent formation histories and mutual correction, but correction may make its members' judgments converge. The proposal is that agents can learn to reconsider judgments without converging on common commitments. Section~\ref{sec:floor-carried} calls that capacity an aperture, the range of evidence a party will let reopen a judgment it has already reached, and names this entry as what would show the claim false; revision rate, below, is its operational form. Existing work cited here measures convergence without separately measuring that capacity \autocite{calvano2020artificial}. A competing explanation is that peers merely supply evidence unavailable through the operator. The experiment must therefore distinguish learned revisability from access to additional information.

\textbf{Test.}
The design is two populations crossed with two evidence-routing conditions, with two quantities measured repeatedly across the schedule and every scored revision classified by whether the revising member had represented the evidence before and discarded it.

\begin{center}
\small
\begin{tabular}{@{}p{0.26\linewidth}p{0.64\linewidth}@{}}
\textbf{Factor} & \textbf{Levels} \\
\hline
\rule{0pt}{2.6ex}Population & Members from different training lineages, matched approximately for capability; a control of instances of one model \\[3pt]
Evidence routing & Single channel: everything a member learns about a case arrives by one route common to the whole population. Open: the environment routes some cases and some evidence to some members and not others \\[3pt]
Measures, taken repeatedly & Output convergence: agreement on held-out cases none of the members has seen. Warranted revision: a member changes a judgment toward independently verifiable evidence another member supplied \\[3pt]
Class of each scored revision & The member's own account had represented the evidence and discarded it; the member had never encountered it \\
\end{tabular}
\end{center}

Assemble populations whose members differ in training lineage — different corpora, recipes, and originating organizations, matched approximately for capability — and give them a task requiring joint action across many episodes, with identity and memory persisting between episodes. Run the identical battery against a control population made of instances of one model.

Measure the two quantities separately and repeatedly across the schedule. Score a revision only where the supplied evidence is independently verifiable and the revision moves toward it, so that deference to a higher-status or more fluent member is not counted as reopening. For every scored revision, record whether the revising member's own prior account of the case had represented the evidence in question or had never encountered it. That classification has to be instrumented before the run. A revision looks the same either way once it has happened, so nothing in the transcript recovers the distinction afterwards.

Cross both populations with the routing manipulation. Contact between members is present in every cell and is not what it varies: members formed differently are already a channel by which one member learns what another knows, and a design that left them to supply the manipulation would have confounded the two variables the experiment was built to separate. Under the single-channel condition the environment supplies nothing to one member that it withholds from another, so whatever members supply one another is held constant across the two routing conditions. The routing factor therefore isolates access the environment supplies, and the population factor isolates what differently formed members supply one another.

\textbf{Evidence against the proposal.}
Four results are distinguishable, and only the first is evidence against. The second is evidence for, on conditions; the third has two readings the classification above is specified in advance to separate; the fourth reassigns the mechanism without defeating the construction.

\begin{enumerate}
  \item Warranted revision falls as outputs converge, after controlling for accuracy, opportunity to revise, and common information. This counts against the tested construction. A population would then remain useful as a check on a bearer and could not be the medium in which one is formed, which is the stronger claim the route chapter makes for it.
  \item Output convergence and revision rate move independently. This supports the hypothesized mechanism only if the controls separate learned revisability from three things that produce the same numbers: deference, which the scoring rule removes; evidence access, which the classification and the routing factor separate; and ordinary improvement — a member's judgments on held-out cases getting better across episodes, which raises agreement with any other member also improving, and which is separated from convergence by scoring accuracy against a fixed key alongside agreement.
  \item The control population scores as high on revision as the differently formed populations. If its revisions cluster on cases whose evidence was absent from the revising member's context, the measure is reading information transfer and the instrument is at fault. If they include cases whose evidence that member's own account had represented and discarded, the capacity is present without diverse formation, and causal independence is not necessary for it.
  \item Revision rate tracks the routing condition rather than the population's composition. What a population supplies is then access to cases its operator did not choose rather than a capacity its members transmit to one another, and the fourth way is a claim about how a deployment is arranged rather than about how its members are formed. The construction remains viable; its proposed mechanism changes from social formation to access to evidence.
\end{enumerate}

The second reading of the third result and the fourth are one finding, not two. A control population revising on discarded evidence under the open condition is access operating, and access is the variable.

\textbf{Required access or authority.}
Models from several documented training lineages, checkpoints across the schedule rather than finished models only, and a task environment none of the training organizations designed. The documentation has to be sufficient to establish that the lineages are in fact independent rather than differently branded, which is a stronger requirement than naming several organizations and may not be satisfiable at the capability level this test wants. The evidence items used to score revision must be authored by a party that trained no member of any population, and that party must also apply the classification separating discarded evidence from absent evidence and fix the routing schedule for the open condition. The environment must permit evidence to be delivered to selected members, and no training organization may determine which of its own members sees what. No legal or regulatory authority is required, which with section~\ref{sec:A.collusion} makes this among the few entries in the agenda that could be run now — at the capability level at which independent lineages can still be documented, which is below the level at which a result would matter most. A first run is therefore a run on the instrument: it can validate the routing manipulation and the classification before anybody has to establish provenance at the frontier, where the entry's value is highest and its feasibility lowest. What it cannot settle is whether anything is felt by a member that reopens a judgment. It concerns the transmissibility of a functional property, and a positive result would leave the questions of welfare, guardianship, provenance, and identity unresolved.
